\section{Autoencoder Framework}

Autoencoders were used to learn compact representations of preprocessed
resting-state fMRI BOLD signals while reconstructing the original data.
As described in Section~\ref{sec:bold_representation}, each session is
represented by a parcellated BOLD time series,
\[
\mathbf{X} \in \mathbb{R}^{T \times R},
\]
where $T$ is the number of timepoints and $R$ is the number of brain
regions. The models operate directly on this representation, rather than
on connectivity matrices or other features derived from it.

An autoencoder has an encoder and a decoder. The encoder $f_{\theta}$
maps the input to a lower-dimensional latent representation,
\[
\mathbf{Z} = f_{\theta}(\mathbf{X}),
\]
where $\mathbf{Z}$ is the compressed representation. The decoder
$g_{\phi}$ maps this representation back to the input space,
\[
\hat{\mathbf{X}} = g_{\phi}(\mathbf{Z}),
\]
where $\hat{\mathbf{X}}$ is the reconstructed BOLD time series. The
complete mapping is
\[
\hat{\mathbf{X}} = g_{\phi}\left(f_{\theta}(\mathbf{X})\right).
\]

The latent representation is an information bottleneck. The models
compress the regional dimension from $R$ to $L$, where $L<R$, while
preserving the temporal dimension. Thus,
\[
\mathbf{X} \in \mathbb{R}^{T \times R}
\longrightarrow
\mathbf{Z} \in \mathbb{R}^{T \times L}
\longrightarrow
\hat{\mathbf{X}} \in \mathbb{R}^{T \times R}.
\]

Each latent feature therefore has a time-dependent trajectory across the scan. The representation describes regional BOLD activity at every timepoint in a lower-dimensional space.
This framework was used for linear, deep fully connected, and convolutional autoencoders. The architectures differ in their encoder and decoder transformations but share the same time-preserving latent format and the common objective of reconstructing the input BOLD signal.

\subsection{Linear Autoencoder}
\thesisfigure{architectures/linearAE}{Linear autoencoder architecture diagram}{fig:lae_diagram}{0.8\textwidth}

The linear autoencoder is the simplest architecture considered. As shown
in Figure~\ref{fig:lae_diagram}, it has one linear encoder layer and one
linear decoder layer. Given an input BOLD time series
\[
\mathbf{X}\in\mathbb{R}^{T\times R},
\]
the encoder maps regional activity at each timepoint from $R$ dimensions
to $L$ dimensions,
\[
\mathbf{X}\in\mathbb{R}^{T\times R}
\longrightarrow
\mathbf{Z}\in\mathbb{R}^{T\times L}.
\]

For each time point $t$, the encoding operation is given by
\[
\mathbf{z}_t = \mathbf{W}_e\mathbf{x}_t + \mathbf{b}_e,
\]
where $\mathbf{x}_t\in\mathbb{R}^{R}$ is regional BOLD activity at
timepoint $t$, $\mathbf{W}_e\in\mathbb{R}^{L\times R}$ is the encoder
weight matrix, and $\mathbf{b}_e\in\mathbb{R}^{L}$ is its bias vector.

The same encoder parameters are applied independently at every
timepoint. This preserves the temporal dimension and makes the number of
trainable parameters independent of $T$.

The decoder maps the latent representation back to the original regional dimensionality,
\[
\hat{\mathbf{x}}_t = \mathbf{W}_d\mathbf{z}_t + \mathbf{b}_d,
\]
where $\mathbf{W}_d\in\mathbb{R}^{R\times L}$ and
$\mathbf{b}_d\in\mathbb{R}^{R}$. Applying the decoder at all timepoints
restores the original dimensionality,
\[
\mathbf{Z}\in\mathbb{R}^{T\times L}
\longrightarrow
\hat{\mathbf{X}}\in\mathbb{R}^{T\times R}.
\]

No nonlinear activation is used. The model therefore learns a
low-dimensional linear representation of regional BOLD activity. Because
each timepoint is processed independently, it does not explicitly model
temporal interactions.

For $R$ input regions and $L$ latent dimensions, the encoder has
$RL+L$ trainable parameters and the decoder has $LR+R$. The total is
$2RL+L+R$.

\subsection{Deep Autoencoder}
\thesisfigure{architectures/deepAE}{Deep autoencoder architecture diagram}{fig:dae_diagram}{\textwidth}

The deep autoencoder extends the linear model with one or two nonlinear
hidden layers before the latent bottleneck. The decoder has the mirrored
structure. Figure~\ref{fig:dae_diagram} shows the evaluated
configurations.

For a single hidden layer, the architecture is
\[
(T,R)\rightarrow(T,H_1)\rightarrow(T,L)\rightarrow(T,H_1)\rightarrow(T,R),
\]
while the two-hidden-layer configuration is
\[
(T,R)\rightarrow(T,H_1)\rightarrow(T,H_2)\rightarrow(T,L)
\rightarrow(T,H_2)\rightarrow(T,H_1)\rightarrow(T,R).
\]

The hidden dimensions $H_1$ and $H_2$ determine the width of the
nonlinear mappings, while $L$ determines the bottleneck size. Hidden
linear layers use GELU activations. The latent projection and final
reconstruction layer remain linear.

As in the linear autoencoder, the same transformations are applied
independently at each timepoint. The temporal dimension is preserved,
but dependencies between neighbouring timepoints are not modelled
explicitly.

The composition of linear transformations and GELU activations allows the
encoder and decoder to learn nonlinear mappings between the regional and
latent representations, increasing representational capacity relative to
the linear autoencoder. Each encoder hidden layer has a corresponding decoder layer.

The experiments used one or two hidden layers to limit model complexity
and computational cost. The latent and hidden-layer dimensions were
varied as architectural hyperparameters.




\subsection{Convolutional Autoencoder}

\subsubsection{Two-Dimensional Convolutional Autoencoder}
\thesisfigure{architectures/conv2dAE}{Conceptual 2D convolutional autoencoder architecture diagram}{fig:conv2dae_diagram}{\textwidth}

The two-dimensional convolutional autoencoder applies convolutions
across the temporal and regional dimensions of the BOLD signal. Its
conceptual architecture is shown in Figure~\ref{fig:conv2dae_diagram}.

The model was motivated by convolutional approaches to neuroimaging
representation learning. It was designed to test whether local features
across time and regional parcel indices improve the learned
representation while retaining a time-preserving latent space.

The architecture preserves the time dimension while reducing the regional
dimension from $R$ to $L$. It therefore produces
$\mathbf{Z}\in\mathbb{R}^{T\times L}$ rather than compressing the whole
scan into one feature map.

The input BOLD signal is represented as
\[
\mathbf{X}\in\mathbb{R}^{T\times R}.
\]
For two-dimensional convolutions, a singleton channel dimension is added,
\[
(B,T,R)\rightarrow(B,1,T,R),
\]
where $B$ is the batch size.

The encoder can contain zero, one, or two hidden convolutional layers.
Each hidden layer applies a two-dimensional convolution and GELU
activation,
\[
\mathbf{H}^{(i)}
=
\operatorname{GELU}
\left(
\operatorname{Conv2D}
\left(
\mathbf{H}^{(i-1)}
\right)
\right).
\]
The kernels extend across time and regions, so a representation can use
neighbouring temporal observations. The temporal stride is one, which
preserves the number of timepoints.

For consistency across architectures, $H_1$ and $H_2$ denote hidden-layer
sizes. In the fully connected architecture, they represent hidden feature
dimensions. In the convolutional architectures, they represent the
number of feature channels. Thus, the 2D hidden convolutional layers
produce $H_1$ and, where used, $H_2$ channels, while $L$ denotes the
width of the final bottleneck.

After the optional hidden layers, a depthwise convolution reduces the
regional dimension. Its kernel and stride operate only along the
regional axis; both are one along the temporal axis. Conceptually,
\[
(H,T,R')
\rightarrow
(H,T,L),
\]
where $R'$ is the regional width after hidden convolutions and $H$ is the
number of channels. Depthwise convolution processes each channel
independently.

A subsequent $1\times1$ convolution combines the feature channels into
one latent channel without changing the temporal or regional dimensions.
The resulting representation is
\[
\mathbf{Z}\in\mathbb{R}^{T\times L}.
\]
This gives one $L$-dimensional vector at each timepoint. The model
reshapes tensors between the convolutional layout and
$\mathbf{Z}\in\mathbb{R}^{T\times L}$ as required.

The decoder reverses this process. A $1\times1$ convolution expands the
latent representation to the required number of channels, and a
depthwise transposed convolution restores regional width. Hidden
convolutional layers are reversed in order. Intermediate decoder layers
use GELU activations and the final layer is linear. The output is
reshaped to
\[
\hat{\mathbf{X}}\in\mathbb{R}^{T\times R}.
\]

Unlike the fully connected models, this architecture uses local temporal
context without reducing temporal resolution. Along the regional axis,
convolution uses neighbouring Schaefer parcel indices, which do not
necessarily represent anatomically adjacent regions. The model therefore
learns local temporal--regional-index patterns rather than explicit
anatomical neighbourhoods.

The experiments considered zero, one, or two hidden convolutional
layers. The latent dimensionality, number of hidden layers, and channel
dimensions were architectural hyperparameters.

\subsubsection{One-Dimensional Regional Convolutional Autoencoder}
\thesisfigure{architectures/convAE}{Conceptual 1D convolutional autoencoder architecture diagram}{fig:convae_diagram}{\textwidth}

The one-dimensional convolutional autoencoder is an intermediate model
between timepoint-wise fully connected autoencoders and the
two-dimensional convolutional model. It tests regional convolution
without temporal convolution. It was introduced to determine whether
regional convolutional feature extraction improves representation
learning independently of temporal convolution.
Figure~\ref{fig:convae_diagram} shows the conceptual architecture.

As with the linear and deep models, each timepoint is processed
independently with shared encoder and decoder parameters. The input
tensor
\[
\mathbf{X}\in\mathbb{R}^{B\times T\times R}
\]
is internally reshaped to
\[
(BT,1,R),
\]
so that each timepoint is an independent one-dimensional regional signal.
The convolutions operate only along the regional axis. After encoding,
the representation is reshaped to
\[
\mathbf{Z}\in\mathbb{R}^{B\times T\times L},
\]
which preserves one $L$-dimensional vector for every timepoint.

The encoder can contain zero, one, or two hidden convolutional layers.
Each hidden layer applies a one-dimensional convolution and GELU
activation,
\[
\mathbf{H}^{(i)}
=
\operatorname{GELU}
\left(
\operatorname{Conv1D}
\left(
\mathbf{H}^{(i-1)}
\right)
\right).
\]
The hidden-layer sizes $H_1$ and $H_2$ denote the number of channels in
their respective convolutional layers. Together with the number of
hidden layers, they control feature-extractor capacity. The latent
dimensionality $L$ determines bottleneck size.

After optional hidden feature extraction, a depthwise one-dimensional
convolution reduces regional width from $R$ to $L$ while processing each
channel independently,
\[
(BT,H,R)\rightarrow(BT,H,L),
\]
where $H$ is the number of channels. The reduction-layer kernel and
stride are selected to give exactly $L$ output positions.

A $1\times1$ convolution then combines the feature channels without
changing regional width,
\[
(BT,H,L)\rightarrow(BT,1,L).
\]
After removing the channel dimension and restoring batch and temporal
dimensions, this gives
\[
\mathbf{Z}\in\mathbb{R}^{B\times T\times L}.
\]
The latent representation at timepoint $t$ therefore depends only on the
regional BOLD values at that timepoint,
\[
\mathbf{z}_t=f_{\theta}(\mathbf{x}_t),
\]
and not on neighbouring observations.

The decoder reverses this process. A $1\times1$ convolution expands the
latent representation to the required number of channels, and a
depthwise transposed convolution restores the regional dimension from
$L$ to $R$. Hidden layers are reversed to return to one output channel.
Intermediate layers use GELU activations and the final layer is linear.
The reconstructed tensor is reshaped to
\[
\hat{\mathbf{X}}\in\mathbb{R}^{B\times T\times R}.
\]

This architecture retains the timepoint-wise processing of the linear and
deep models but replaces fully connected layers with regional
convolutions. Its filters share parameters across regional positions,
giving locally constrained feature extraction rather than the global
mappings used by fully connected layers.

Unlike the two-dimensional model, it does not learn convolutional
temporal features. The distinction is
\[
\text{1D ConvAE: regional convolution only}
\]
and
\[
\text{2D ConvAE: joint temporal--regional convolution}.
\]
The one-dimensional model therefore separates the effect of regional
convolution from temporal convolution.

As convolution follows the Schaefer parcel ordering, neighbouring
positions are not necessarily anatomically adjacent. The model learns
local regional-index patterns, not explicit anatomical neighbourhoods.

The experiments considered zero, one, or two hidden convolutional
layers. The latent dimensionality, number of layers, and $H_1$ and $H_2$
were architectural hyperparameters.

\subsection{Architecture Summary}

The principal characteristics of the deterministic autoencoder architectures
considered in this work are summarised in Table~\ref{tab:ae_architecture_summary}.
All architectures preserve a latent representation of size $T\times L$, but
differ in how information is transformed across the regional and temporal
dimensions. The linear and deep autoencoders process each time point
independently using fully connected transformations, the one-dimensional
convolutional autoencoder introduces convolutional processing along the
regional dimension only, and the two-dimensional convolutional autoencoder
additionally incorporates local temporal context.

\begin{table}[ht]
    \centering
    \caption{Summary of the deterministic autoencoder architectures.}
    \label{tab:ae_architecture_summary}

    \begin{tabular}{
        >{\raggedright\arraybackslash}p{0.17\textwidth}
        >{\raggedright\arraybackslash}p{0.19\textwidth}
        >{\raggedright\arraybackslash}p{0.25\textwidth}
        >{\raggedright\arraybackslash}p{0.15\textwidth}
        >{\raggedright\arraybackslash}p{0.11\textwidth}
    }
        \toprule
        \textbf{Architecture} &
        \textbf{Temporal modelling} &
        \textbf{Regional mapping} &
        \textbf{Hidden layers} &
        \textbf{Latent} \\
        \midrule

        Linear AE &
        None &
        Linear, fully connected &
        None &
        $T\times L$ \\

        Deep AE &
        None &
        Nonlinear, fully connected &
        1--2 &
        $T\times L$ \\

        1D Regional ConvAE &
        None &
        1D regional convolution &
        0--2 &
        $T\times L$ \\

        2D ConvAE &
        Local temporal context &
        2D temporal--regional convolution &
        0--2 &
        $T\times L$ \\

        \bottomrule
    \end{tabular}
\end{table}


These architectures therefore form a progression from globally connected,
timepoint-wise mappings to convolutional representations with increasing
use of local structure. The following section introduces the variational
extension applied to each architecture while retaining the same underlying
encoder-decoder configuration and time-preserving latent dimensionality.




\subsection{Variational Autoencoder Extension}
\thesisfigure{architectures/variationalAE}{Conceptual variational extension applied to the autoencoder architectures}{fig:variational_ae_diagram}{0.9\textwidth}

After the architectural search, selected deterministic autoencoders were
extended with the variational autoencoder formulation of Kingma and
Welling~\cite{KingmaWelling2014}. This tests whether a structured
probabilistic latent space improves representation organisation and
generalisation while retaining reconstruction quality. Figure~\ref{fig:variational_ae_diagram}
shows this modification.

In deterministic autoencoders, the encoder directly produces
\[
\mathbf{Z}=f_{\theta}(\mathbf{X}).
\]
In the variational formulation, it instead parameterises an approximate
posterior distribution,
\[
q_{\theta}(\mathbf{Z}\mid\mathbf{X})
=
\mathcal{N}
\left(
\boldsymbol{\mu},
\operatorname{diag}(\boldsymbol{\sigma}^{2})
\right),
\]
where $\boldsymbol{\mu}$ and $\boldsymbol{\sigma}$ are the latent mean
and standard deviation. The latent representation is sampled from this
distribution.

The models predict $\boldsymbol{\mu}$ and log-variance
$\log\boldsymbol{\sigma}^{2}$ rather than standard deviation directly.
The standard deviation is obtained as
\[
\boldsymbol{\sigma}
=
\exp\left(
\frac{1}{2}\log\boldsymbol{\sigma}^{2}
\right).
\]
This keeps the network output unconstrained while ensuring a positive
standard deviation. Log-variance is bounded for numerical stability.

Latent samples use the reparameterisation trick,
\[
\boldsymbol{\epsilon}
\sim
\mathcal{N}(\mathbf{0},\mathbf{I}),
\]
\[
\mathbf{Z}
=
\boldsymbol{\mu}
+
\boldsymbol{\sigma}\odot\boldsymbol{\epsilon},
\]
where $\odot$ denotes element-wise multiplication. This separates
sampling from the learnable parameters, allowing gradients to propagate
through $\boldsymbol{\mu}$ and $\boldsymbol{\sigma}$.

The variational extension does not change the latent representation
shape. For the time-preserving models,
\[
\boldsymbol{\mu},
\log\boldsymbol{\sigma}^{2},
\mathbf{Z}
\in
\mathbb{R}^{T\times L},
\]
so each input timepoint retains an $L$-dimensional latent vector. The
extension changes the parameterisation and regularisation of the
bottleneck, not its size or temporal structure.

The approximate posterior is regularised towards a standard normal
prior,
\[
p(\mathbf{Z})
=
\mathcal{N}(\mathbf{0},\mathbf{I}).
\]
using the Kullback--Leibler (KL) divergence,
\[
\mathcal{L}_{\mathrm{KL}}
=
D_{\mathrm{KL}}
\left(
q_{\theta}(\mathbf{Z}\mid\mathbf{X})
\parallel
p(\mathbf{Z})
\right).
\]
For a diagonal Gaussian posterior and standard normal prior, the KL term
is
\[
\mathcal{L}_{\mathrm{KL}}
=
-\frac{1}{2}
\sum_j
\left(
1+\log\sigma_j^2-\mu_j^2-\sigma_j^2
\right),
\]
The implementation normalises this term over the latent dimensions.

The training objective adds this term to reconstruction loss,
\[
\mathcal{L}
=
\mathcal{L}_{\mathrm{recon}}
+
\beta\mathcal{L}_{\mathrm{KL}},
\]
where $\beta$ controls variational regularisation. Smaller values favour
reconstruction fidelity, while larger values constrain the posterior more
strongly towards the standard normal prior.

During training, latent values are sampled with this procedure. During
evaluation, the posterior mean $\boldsymbol{\mu}$ is used instead,
giving a deterministic representation and reconstruction.
